\subsection{Interfaces}

\subsubsection{Regla: prefijo \texttt{I}}

\textbf{Regla:} Toda interfaz se escribe en \texttt{PascalCase} con prefijo \texttt{I}, seguido de un sustantivo, frase nominal o adjetivo que describa la capacidad o rol (Cwalina \& Abrams, 2008).

\textbf{Código correcto:}
\begin{lstlisting}[style=csharpstyle]
public interface IMovable
{
    void Move(Tile destination);
}
\end{lstlisting}

\textbf{Código incorrecto:}
\begin{lstlisting}[style=csharpstyle]
public interface Movable
{
    void Move(Tile destination);
}
\end{lstlisting}

\subsubsection{Regla: correspondencia entre interfaz e implementación estándar}

\textbf{Regla:} Cuando una clase es la implementación estándar de una interfaz, su nombre coincide exactamente con el de la interfaz, sin el prefijo \texttt{I} (Cwalina \& Abrams, 2008).

\textbf{Código correcto:}
\begin{lstlisting}[style=csharpstyle]
public interface IDiceRoller
{
    int Roll();
}

public sealed class DiceRoller : IDiceRoller
{
    public int Roll()
    {
        return Random.Shared.Next(1, 7);
    }
}
\end{lstlisting}

\textbf{Código incorrecto:}
\begin{lstlisting}[style=csharpstyle]
public interface IDiceRoller
{
    int Roll();
}

public sealed class RandomDiceGenerator : IDiceRoller
{
    public int Roll()
    {
        return Random.Shared.Next(1, 7);
    }
}
\end{lstlisting}
